\documentclass[10pt,a4paper]{amsart}
\usepackage{amsmath,amssymb,amsthm,graphicx}
\usepackage{verbatim}
\newenvironment{reference}{\comment}{\endcomment}
%\newenvironment{reference}{}{}
\newenvironment{slogan}{\comment}{\endcomment}
\newenvironment{history}{\comment}{\endcomment}

% For commutative diagrams we use Xy-pic
\usepackage[all]{xy}

% We use 2cell for 2-commutative diagrams.
\xyoption{2cell}
% Theorem environments.
%
\theoremstyle{plain}
\newtheorem{theorem}[subsection]{Theorem}
\newtheorem{proposition}[subsection]{Proposition}
\newtheorem{lemma}[subsection]{Lemma}

\theoremstyle{definition}
\newtheorem{definition}[subsection]{Definition}
\newtheorem{example}[subsection]{Example}
\newtheorem{exercise}[subsection]{Exercise}
\newtheorem{situation}[subsection]{Situation}

\theoremstyle{remark}
\newtheorem{remark}[subsection]{Remark}
\newtheorem{remarks}[subsection]{Remarks}

\numberwithin{equation}{subsection}

% Macros
%
\def\lim{\mathop{\mathrm{lim}}\nolimits}
\def\colim{\mathop{\mathrm{colim}}\nolimits}
\def\Spec{\mathop{\mathrm{Spec}}}
\def\Hom{\mathop{\mathrm{Hom}}\nolimits}
\def\Ext{\mathop{\mathrm{Ext}}\nolimits}
\def\SheafHom{\mathop{\mathcal{H}\!\mathit{om}}\nolimits}
\def\SheafExt{\mathop{\mathcal{E}\!\mathit{xt}}\nolimits}
\def\Sch{\mathit{Sch}}
\def\Mor{\mathop{\mathrm{Mor}}\nolimits}
\def\Ob{\mathop{\mathrm{Ob}}\nolimits}
\def\Sh{\mathop{\mathit{Sh}}\nolimits}
\def\NL{\mathop{N\!L}\nolimits}
\def\CH{\mathop{\mathrm{CH}}\nolimits}
\def\proetale{{pro\text{-}\acute{e}tale}}
\def\etale{{\acute{e}tale}}
\def\QCoh{\mathit{QCoh}}
\def\Ker{\mathop{\mathrm{Ker}}}
\def\Im{\mathop{\mathrm{Im}}}
\def\Coker{\mathop{\mathrm{Coker}}}
\def\Coim{\mathop{\mathrm{Coim}}}

% Boxtimes
%
\DeclareMathSymbol{\boxtimes}{\mathbin}{AMSa}{"02}

%
% Macros for moduli stacks/spaces
%
\def\QCohstack{\mathcal{QC}\!\mathit{oh}}
\def\Cohstack{\mathcal{C}\!\mathit{oh}}
\def\Spacesstack{\mathcal{S}\!\mathit{paces}}
\def\Quotfunctor{\mathrm{Quot}}
\def\Hilbfunctor{\mathrm{Hilb}}
\def\Curvesstack{\mathcal{C}\!\mathit{urves}}
\def\Polarizedstack{\mathcal{P}\!\mathit{olarized}}
\def\Complexesstack{\mathcal{C}\!\mathit{omplexes}}
% \Pic is the operator that assigns to X its picard group, usage \Pic(X)
% \Picardstack_{X/B} denotes the Picard stack of X over B
% \Picardfunctor_{X/B} denotes the Picard functor of X over B
\def\Pic{\mathop{\mathrm{Pic}}\nolimits}
\def\Picardstack{\mathcal{P}\!\mathit{ic}}
\def\Picardfunctor{\mathrm{Pic}}
\def\Deformationcategory{\mathcal{D}\!\mathit{ef}}
\usepackage{cleveref}
\crefname{originalsection}{original section}{original sections}
\Crefname{originalsection}{Original section}{Original sections}
\makeatletter
\newlabel{algebra-section-netherlander}{{00S0}{source}{}{sourcepointer-1}{}}
\newlabel{algebra-section-netherlander@cref}{{[originalsection][0][]00S0}{[1][0][]source}{}{}{}}
\makeatother
\crefname{originallemma}{original lemma}{original lemmas}
\Crefname{originallemma}{Original lemma}{Original lemmas}
\makeatletter
\newlabel{more-algebra-lemma-silly}{{0G9L}{source}{}{sourcepointer-2}{}}
\newlabel{more-algebra-lemma-silly@cref}{{[originallemma][0][]0G9L}{[1][0][]source}{}{}{}}
\makeatother
\crefname{originallemma}{original lemma}{original lemmas}
\Crefname{originallemma}{Original lemma}{Original lemmas}
\makeatletter
\newlabel{more-algebra-lemma-two-term-surjection-map-zero}{{0G9I}{source}{}{sourcepointer-3}{}}
\newlabel{more-algebra-lemma-two-term-surjection-map-zero@cref}{{[originallemma][0][]0G9I}{[1][0][]source}{}{}{}}
\makeatother
\title{The Jacobian determinant annihilates the naive cotangent complex of a presentation}
\author{The Stacks Project Authors}
\date{stacks-project:restricted.tex:local-snapshot-20261012; GNU FDL 1.2 or later}
\begin{document}
\maketitle
\noindent\textit{Editorial note.} All delivered mathematical text is the author's own, copied from the pinned Stacks Project chapter file restricted.tex; no statement or proof step was rewritten, reordered or removed.\par\medskip
\noindent\textit{Editorial note.} The counted proof is the author's complete argument: the proof transfers the vanishing of multiplication by the Jacobian determinant from the matrix complex K through the two-term presentation complex L up to Ext\textasciicircum{}1 of the full naive cotangent complex via a distinguished triangle\par\medskip
\noindent\textit{Editorial note.} The label lemma-presentation-rig-smooth names the counted statement, and the proof environment immediately following it in the source is delivered verbatim as its proof.\par\medskip
\noindent\textit{Editorial note.} The counted statement is carried by the author's own lemma environment in the chapter source.\par\medskip
\noindent\textit{Allowed standard prerequisites (editorial).}\begin{itemize}
\item the naive cotangent complex NL of a ring map and Ext groups over D(B)
\item the two-term complex representing NL of a quotient presentation (Algebra, Section algebra-section-netherlander)
\item the More on Algebra lemmas: the determinant of the Jacobian matrix acts trivially on the displayed matrix complex (more-algebra-lemma-silly) and a surjection of two-term complexes transfers vanishing of multiplication (more-algebra-lemma-two-term-surjection-map-zero)
\item distinguished triangles in the derived category and the long exact Ext sequence
\end{itemize}\medskip
\ifdefined\hypertarget\hypertarget{sourcepointer-1}{}\fi
\noindent\textit{Source reference.} Original section 00S0: sources/stacks-chow/tags.txt, source lines 1022--1022. Prerequisite pointer invoked at proof line(s) 1617; it names a standard fact of a sibling chapter that the counted argument uses, and every step of the counted proof built on it is retained verbatim in the delivered body.\par\medskip
\ifdefined\hypertarget\hypertarget{sourcepointer-2}{}\fi
\noindent\textit{Source reference.} Original lemma 0G9L: sources/stacks-chow/tags.txt, source lines 20306--20306. Prerequisite pointer invoked at proof line(s) 1613; it names a standard fact of a sibling chapter that the counted argument uses, and every step of the counted proof built on it is retained verbatim in the delivered body.\par\medskip
\ifdefined\hypertarget\hypertarget{sourcepointer-3}{}\fi
\noindent\textit{Source reference.} Original lemma 0G9I: sources/stacks-chow/tags.txt, source lines 20303--20303. Prerequisite pointer invoked at proof line(s) 1627; it names a standard fact of a sibling chapter that the counted argument uses, and every step of the counted proof built on it is retained verbatim in the delivered body.\par\medskip
\begin{lemma}
\label{lemma-presentation-rig-smooth}
Let $A$ be a ring. Let $f_1, \ldots, f_m \in A[x_1, \ldots, x_n]$
and set $B = A[x_1, \ldots, x_n]/(f_1, \ldots, f_m)$. Assume $m \leq n$
and set $g = \det_{1 \leq i, j \leq m}(\partial f_j/\partial x_i)$.
Then
\begin{enumerate}
\item $g$ annihilates $\Ext^1_B(\NL_{B/A}, N)$ for every $B$-module $N$,
\item if $n = m$, then multiplication by $g$ on $\NL_{B/A}$ is $0$ in $D(B)$.
\end{enumerate}
\end{lemma}

\begin{proof}
Let $T$ be the $m \times m$ matrix with entries $\partial f_j/\partial x_i$
for $1 \leq i, j \leq n$. Let $K \in D(B)$ be represented by the complex
$T : B^{\oplus m} \to B^{\oplus m}$ with terms sitting in degrees $-1$ and $0$.
By More on Algebra, Lemmas \ref{more-algebra-lemma-silly}
we have $g : K \to K$ is zero in $D(B)$. Set $J = (f_1, \ldots, f_m)$.
Recall that $\NL_{B/A}$ is homotopy equivalent to
$J/J^2 \to \bigoplus_{i = 1, \ldots, n} B\text{d}x_i$, see
Algebra, Section \ref{algebra-section-netherlander}.
Denote $L$ the complex $J/J^2 \to \bigoplus_{i = 1, \ldots, m} B\text{d}x_i$
to that we have the quotient map $\NL_{B/A} \to L$.
We also have a surjective map of complexes $K \to L$
by sending the $j$th basis element in the term $B^{\oplus m}$
in degree $-1$ to the class of $f_j$ in $J/J^2$. Picture
$$
\NL_{B/A} \to L \leftarrow K
$$
From
More on Algebra, Lemma \ref{more-algebra-lemma-two-term-surjection-map-zero}
we conclude that multiplication by $g$ on $L$ is $0$ in $D(B)$.
On the other hand, the distinguished triangle
$B^{\oplus n - m}[0] \to \NL_{B/A} \to L$
shows that $\Ext^1_B(L, N) \to \Ext^1_B(\NL_{B/A}, N)$ is surjective
for every $B$-module $N$ and hence annihilated by $g$.
This proves part (1).
If $n = m$ then $\NL_{B/A} = L$ and we see that (2) holds.
\end{proof}

\end{document}
